\def\examid{26101200002}
\begin{document}
\begin{enumerate}
\item \multipart{1}{8}
The data set \code{bwght} contains data on births in the US from the 1988 National Health Interview Survey. You can load the data into R using the \code{wooldridge} library using:

\medskip
\code{library(wooldridge)}\\
\code{dat = wooldridge::bwght}
\medskip

The main outcome of interest is the variable \textit{bwght}, which is the birth weight of the child in ounces.\\
What is the mean of the variable \textit{bwght} in the dataset?
\begin{opts}
\item 120.0000
\item \ca{118.6996}
\item 116.9325
\item 119.5121
\item None of the above
\end{opts}

\item How many missing values or NA's are there for the variable \textit{fatheduc} in the dataset?
\begin{opts}
\item 0
\item 1
\item \ca{196}
\item 197
\item None of the above
\end{opts}

\item Consider the following linear regression model.
\[\begin{aligned}
bwght = \beta_0 + \beta_1(cigs) + \beta_2(faminc) + \beta_3(faminc^2) + \beta_4(parity) &+ \\
\beta_5(motheduc) + \beta_6(fatheduc) + \beta_7(male) + \beta_8(white) &+ \varepsilon
\end{aligned}\]
The model includes the number of cigarettes smoked per day by the mother during pregnancy (\textit{cigs}), family income and family income squared, and the birth order of the child (\textit{parity}). The model also controls for the education of the mother and father. Finally, (\textit{male, white}), indicator variables for the sex and race of the child are included in the model.\\
Estimate the model above with OLS. What is the estimated coefficient on \textit{cigs}?
\begin{opts}
\item $-0.4634$
\item \ca{$-0.5950$}
\item $-0.5137$
\item $-0.6211$
\item $-0.4772$
\end{opts}

\item What is the sum of squared errors of the model?
\begin{opts}
\item 461992.3
\item 13340.2
\item 390114.8
\item 511023.7
\item \ca{None of the above}
\end{opts}

\item Suppose we wanted to create a confidence interval for one of the parameters of the model. If we wanted to create a 99\% confidence interval, what critical value would we use?
\begin{opts}
\item 2.5758
\item 2.3295
\item \ca{2.5800}
\item 1.9620
\item None of the above
\end{opts}

\item Consider the null hypothesis that $H_0:\beta_{male}=0$. What is the strongest conclusion we can draw about this null hypothesis?
\begin{opts}
\item \ca{Reject the null at the 0.1\% level}
\item Reject the null at the 1\% level
\item Reject the null at the 5\% level
\item Reject the null at the 10\% level
\item Fail to reject the null at the 10\% level
\end{opts}

\item Consider the null hypothesis that $H_0:\beta_{cigs}=-0.5$. What is the p-value associated with this hypothesis test?
\begin{opts}
\item 0.1935
\item \ca{0.3871}
\item 0.6129
\item 0.0000
\item None of the above
\end{opts}

\item What is the covariance between the estimate of the parameter on \textit{(Intercept)} and the parameter on \textit{white}?
\begin{opts}
\item $-0.9241$
\item 2.6122
\item \ca{$-1.6775$}
\item 1.6775
\item $-0.0017$
\end{opts}

\item \multipart{9}{15}
The data set \code{k401ksubs} contains data on workers from the 1991 Survey of Income and Program Participation. You can load the data into R using the \code{wooldridge} library using:

\medskip
\code{library(wooldridge)}\\
\code{dat = wooldridge::k401ksubs}
\medskip

The variable of interest is \textit{e401k} which indicates if the worker is eligible for a 401(k) plan through their employer. We would like to study what are the main determinants of 401(k) eligibility using a linear probability model. To answer this question you will need to add a new variable to the data set called \textit{young} which should equal 1 if the worker is younger than 35 (i.e., \textit{age} $<35$) and zero otherwise.\\
What percent of the population are \textit{young}?
\begin{opts}
\item 28.42\%
\item 68.31\%
\item 35.11\%
\item \ca{31.69\%}
\item None of the above
\end{opts}

\item Estimate a linear probability model on 401(k) eligibility \textit{e401k}.\\
Include in your model controls for \textit{young}, annual family income in \$1000s (\textit{inc}), and indicators for whether the worker is male (\textit{male}) and married (\textit{marr}). Your model should have 5 parameters including the intercept.\\
What is the R-squared from the estimated model?
\begin{opts}
\item 0.07355
\item \ca{0.07394}
\item 0.06112
\item 0.08227
\item 0.05930
\end{opts}

\item What is the correct interpretation of the parameter on \textit{young}? Holding fixed the variables in the model, being younger than 35 changes the probability of being eligible for a 401(k) by \dots
\begin{opts}
\item $-4.0627$ percent
\item $-0.0406$ percentage points
\item \ca{$-4.0627$ percentage points}
\item $-0.4063$ percentage points
\item None of the above
\end{opts}

\item What is the estimate of the upper bound of the 99\% confidence interval for the parameter on \textit{marr}?
\begin{opts}
\item $-0.0015$
\item $-0.0543$
\item 0.0092
\item \ca{0.0057}
\item None of the above
\end{opts}

\item Consider the null hypothesis that $H_0:\beta_{marr}=0$. What is the strongest conclusion we can draw about this null hypothesis?
\begin{opts}
\item Reject the null at the 0.1\% level
\item Reject the null at the 1\% level
\item \ca{Reject the null at the 5\% level}
\item Reject the null at the 10\% level
\item Fail to reject the null at the 10\% level
\end{opts}

\item Re-estimate the model adding family size (\textit{fsize}) as an additional control. What is the percentage change in the parameter on \textit{marr} when we control for family size?
\begin{opts}
\item \ca{$-41.27\%$}
\item 41.27\%
\item $-29.12\%$
\item $-70.27\%$
\item None of the above
\end{opts}

\item Re-estimate the model adding family size (\textit{fsize}) as an additional control. What is the variance inflation factor for the parameter on \textit{marr} when we control for family size? (note: if you are calculating the variance inflation factor from the formula you need to ensure that the same observations are used in the different regressions)
\begin{opts}
\item 1.5284
\item 1.1983
\item \ca{1.8029}
\item 0.5547
\item None of the above
\end{opts}
\end{enumerate}

\keyheader
\begin{keytable}
\krow{1}{b}{118.6996}{\code{mean(dat\$bwght)}; 120 is the median}
\krow{2}{c}{196}{\code{sum(is.na(dat\$fatheduc))}}
\krow{3}{b}{$-0.5950$}{lm drops 197 rows (196 \textit{fatheduc} + 1 \textit{motheduc}), $n=1191$}
\krow{4}{e}{456161.5}{\code{deviance(m)}; not listed}
\krow{5}{c}{2.5800}{\code{qt(0.995, df.residual(m))}, df $=1191-9=1182$; 2.5758 is \code{qnorm}}
\krow{6}{a}{$p=0.00077$}{$p<0.001$}
\krow{7}{b}{0.3871}{\code{linearHypothesis(m, "cigs = -0.5")}, \code{Pr(>F)}. 0.1935 is one-sided}
\krow{8}{c}{$-1.6775$}{\code{vcov(m)["(Intercept)","white"]}}
\krow{9}{d}{31.69\%}{\code{young = as.numeric(age < 35)}; \code{<=35} gives 35.11\%}
\krow{10}{b}{0.07394}{0.07355 is adjusted $R^2$}
\krow{11}{c}{$-0.0406$}{LPM: $\times100$ = percentage points}
\krow{12}{d}{0.0057}{\code{confint(m1, "marr", level=0.99)}; 0.0092 is the 95\% upper bound on \textit{male}, $-0.0015$ the 95\% upper on \textit{marr}}
\krow{13}{c}{$p=0.0367$}{$0.01<p<0.05$}
\krow{14}{a}{$-41.27\%$}{\code{(new - old) / old * 100} $=(-0.014287+0.024325)/(-0.024325)\times100$}
\krow{15}{c}{1.8029}{\code{car::vif(m2)}; 0.5547 is $1-R_j^2$}
\end{keytable}

\subsection*{R code}
\lstinputlisting{../key_practiceA.r}
\end{document}
